% Put the results and discussion here

\section{RESULTS AND DISCUSSION}

This section presents the findings from the pilot user evaluation of the Digital Farmers' Produce System. The discussion focuses on the respondents' profile, user acceptance results based on the Technology Acceptance Model (TAM), system functionality, overall satisfaction, and insights gained during the conduct of the project. Since the evaluation involved a small purposive sample of intended users, the findings are interpreted as pilot evidence rather than statistically generalizable results.

\subsection{Respondent Profile}

The user testing involved twelve (12) respondents composed of five (5) farmers, four (4) buyers, and three (3) Department of Agriculture/Municipal Agriculture Office (DA/MAO) staff. The respondents represented the three major user roles of the system. This was important because the Digital Farmers' Produce System was designed to support interactions among farmers, buyers, and agricultural officers.

\begin{table}[htbp]
\vspace{2ex}
\centering
\caption{Profile of Respondents by Age Group and Role}
\label{table:respondent-age-role}
\begin{tabular}{p{3.2cm}ccccc}
\hline
\hline
\textbf{Age Group} & \textbf{Farmers} & \textbf{Buyers} & \textbf{DA/MAO} & \textbf{Total f} & \textbf{Total \%} \\
\hline
18--25 years old & 0 & 0 & 0 & 0 & 0.0 \\
26--35 years old & 0 & 0 & 2 & 2 & 16.7 \\
36--45 years old & 1 & 3 & 0 & 4 & 33.3 \\
46 years old and above & 4 & 1 & 1 & 6 & 50.0 \\
\hline
\textbf{Total} & \textbf{5} & \textbf{4} & \textbf{3} & \textbf{12} & \textbf{100.0} \\
\hline
\hline
\end{tabular}
\vspace{1ex}

\begin{flushleft}
\footnotesize{\textit{Note.} Percentages are based on N = 12.}
\end{flushleft}
\vspace{2ex}
\end{table}

As shown in Table \ref{table:respondent-age-role}, half of the respondents were 46 years old and above. No respondent belonged to the 18--25 age group. This profile is relevant because the system is intended for actual agricultural users, many of whom may not be digital natives. Thus, favorable usability results from an older user group provide practical evidence that the simplified interface can support the target community.

\begin{table}[htbp]
\vspace{2ex}
\centering
\caption{Experience in Using Smartphones/Web Applications by Frequency of System Use}
\label{table:experience-frequency}
\begin{tabular}{p{3.8cm}ccccc}
\hline
\hline
\textbf{Frequency of System Use} & \textbf{Beginner} & \textbf{Intermediate} & \textbf{Advanced} & \textbf{Total f} & \textbf{Total \%} \\
\hline
Daily & 0 & 0 & 0 & 0 & 0.0 \\
Weekly & 0 & 0 & 0 & 0 & 0.0 \\
Occasionally & 0 & 0 & 0 & 0 & 0.0 \\
First-time user & 4 & 8 & 0 & 12 & 100.0 \\
\hline
\textbf{Total} & \textbf{4} & \textbf{8} & \textbf{0} & \textbf{12} & \textbf{100.0} \\
\hline
\hline
\end{tabular}
\vspace{1ex}

\begin{flushleft}
\footnotesize{\textit{Note.} All respondents evaluated the prototype as first-time users.}
\end{flushleft}
\vspace{2ex}
\end{table}

Table \ref{table:experience-frequency} shows that all respondents were first-time users of the system. Four respondents reported beginner experience in using smartphones or web applications, while eight reported intermediate experience. This strengthens the practical meaning of the usability ratings because the respondents evaluated a system with which they were not previously familiar. At the same time, the results suggest that sustained adoption will depend on proper onboarding, practice sessions, and continued support from the local agriculture office.

\subsection{Technology Acceptance Results}

The Technology Acceptance Model was used to evaluate user acceptance of the system. The constructs measured were perceived usefulness, perceived ease of use, attitude toward using, and behavioral intention to use. Additional categories were included to evaluate system functionality and overall satisfaction.

\begin{table}[htbp]
\vspace{2ex}
\centering
\caption{Summary of Mean Ratings for System Usage by Role}
\label{table:summary-mean-ratings}
\begin{tabular}{p{4.6cm}ccc}
\hline
\hline
\textbf{Indicator} & \textbf{Farmers M} & \textbf{Buyers M} & \textbf{DA/MAO M} \\
\hline
Perceived Usefulness & 4.45 & 4.81 & 4.58 \\
Perceived Ease of Use & 4.55 & 4.81 & 4.75 \\
Attitude Toward Using & 4.87 & 5.00 & 4.67 \\
Behavioral Intention to Use & 4.40 & 4.50 & 4.78 \\
System Functionality Evaluation & 4.84 & 4.95 & 4.67 \\
Overall Satisfaction & 4.50 & 4.88 & 4.67 \\
\hline
\textbf{Overall Mean} & \textbf{4.60} & \textbf{4.83} & \textbf{4.69} \\
\hline
\hline
\end{tabular}
\vspace{1ex}

\begin{flushleft}
\footnotesize{\textit{Note.} Scale: 1.00--1.80 = Strongly Disagree; 1.81--2.60 = Disagree; 2.61--3.40 = Neutral; 3.41--4.20 = Agree; 4.21--5.00 = Strongly Agree. M = Mean.}
\end{flushleft}
\vspace{2ex}
\end{table}

As shown in Table \ref{table:summary-mean-ratings}, all role groups rated the system positively. Farmers obtained an overall mean of 4.60, buyers obtained 4.83, and DA/MAO staff obtained 4.69. All of these results fall under the interpretation of \textit{Strongly Agree}. This indicates that the intended users strongly accepted the Digital Farmers' Produce System during the pilot evaluation.

For perceived usefulness, farmers recorded an overall mean of 4.45, buyers 4.81, and DA/MAO officers 4.58. These results show that all user groups perceived the system as useful. Buyers gave the highest usefulness rating, suggesting that the product search and farmer contact functions directly addressed their need for timely produce visibility. Farmers also rated the system positively, although some variation was observed in the item related to connecting with buyers and/or farmers. This may reflect differences in confidence, previous selling practices, or expectations about whether digital contact will lead to actual transactions.

Perceived ease of use was also high. Farmers recorded a mean of 4.55, buyers 4.81, and DA/MAO staff 4.75. These results indicate that the respondents found the system easy to learn and navigate. The interface clarity item received very high ratings across groups, suggesting that the system design was understandable even to users with limited digital experience. However, the farmer group recorded a slightly lower result on the item concerning use without assistance. This indicates that some farmers may still need guided use during the early stages of implementation.

Attitude toward using the system was the strongest acceptance dimension. Farmers recorded a mean of 4.87, buyers 5.00, and DA/MAO officers 4.67. These results show that the respondents viewed the system as a good idea for the community and felt comfortable using it. This is important because positive attitude can influence continued use and support for implementation.

Behavioral intention to use was also favorable. Farmers recorded a mean of 4.40, buyers 4.50, and DA/MAO staff 4.78. These results suggest that respondents were willing to use the system once available. The relatively lower farmer score indicates that actual long-term use may depend on continued access, visible buyer participation, and support from the agriculture office.

\subsection{System Functionality and Overall Satisfaction}

The system functionality evaluation received strong ratings from all groups. Farmers gave a mean of 4.84, buyers 4.95, and DA/MAO officers 4.67. These results indicate that the respondents strongly agreed that the product listing, search and filtering, messaging, notification, and system response features functioned properly.

The product listing and messaging features received particularly strong responses. This suggests that the system was able to support its core purpose of helping farmers post produce and communicate with potential buyers. The notification features, including SMS, were also rated positively, which is important for users who may have limited access to smartphones or stable internet connections.

The DA/MAO group's rating for search and filtering was positive but slightly lower than other functionality ratings. This suggests that administrative users may need more refined filters, dashboard summaries, or reporting options in future versions. This finding is important because DA/MAO staff use the system not only for ordinary browsing but also for monitoring, reporting, and decision support.

Overall satisfaction was also favorable. Farmers recorded a mean of 4.50, buyers 4.88, and DA/MAO officers 4.67. These results suggest that the prototype met the immediate needs of the intended users. However, the findings should be interpreted with caution because the sample size was small, all users were first-time users, and long-term adoption was not yet measured.

\subsection{Insights and Lessons Learned}

Several insights were gained during the conduct of the project. First, the development of a localized agricultural platform requires a clear understanding of the actual practices of farmers, buyers, and agricultural officers. The system cannot focus only on buying and selling; it must also support coordination, monitoring, and communication among stakeholders.

Second, simplicity is essential for user acceptance. Since some respondents had only beginner-level experience with smartphones or web applications, the system needed clear labels, direct workflows, and mobile-friendly layouts. The positive ease-of-use ratings suggest that a simplified interface can help reduce barriers to adoption.

Third, the project showed that market linkage depends on participation from both sides. Farmers will benefit more if buyers actively search the system, while buyers will benefit more if farmers regularly update product listings. This means that the local agriculture office has an important role in onboarding users, encouraging participation, and maintaining data quality.

Fourth, the system can support municipal-level agricultural coordination by giving DA/MAO staff better visibility of available produce, farmer participation, and resource requests. This can help improve planning, monitoring, and possible subsidy distribution.

\subsection{Challenges Encountered}

Several challenges were encountered during the project. One challenge was designing the system for users with different levels of digital literacy. Farmers, buyers, and DA/MAO personnel have different expectations and ways of using technology. To address this, the interface was designed to be simple, responsive, and role-based.

Another challenge was ensuring that the system would work properly on mobile devices. Since many intended users are expected to access the system through smartphones, responsive design was necessary. The layout had to be adjusted to ensure that buttons, forms, product listings, and messages remained readable and usable on smaller screens.

A further challenge was maintaining accurate data. The usefulness of the system depends on the accuracy of product listings, quantities, prices, and harvest schedules entered by users. This means that future deployment should include user training and possible verification procedures by the agriculture office.

Connectivity was also considered a limitation. Since some rural users may experience unstable internet access, the system must continue to improve its support for low-connectivity environments. SMS notifications and simplified pages can help address this issue, but further enhancements may still be needed.

\subsection{Implications for Agricultural Coordination}

The findings suggest three major implications for agricultural coordination in Sinacaban. First, localized digital agriculture systems can be accepted by rural users when they are designed around familiar tasks, simple navigation, and mobile-first workflows. Second, market linkage is not only a farmer-buyer activity; institutional visibility is also important because DA/MAO staff need structured information for planning, resource allocation, and monitoring. Third, the usefulness of the system depends on active participation. The system will be more effective if farmers regularly post updated produce information and buyers consistently use the platform to search for available crops.

\subsection{Limitations of the Evaluation}

The evaluation has several limitations. The pilot sample was small and was drawn from intended local users rather than from a random population. Therefore, the results should be treated as pilot evidence. The questionnaire results were summarized using means and standard deviations, which limits more advanced reliability and inferential analysis. Since only aggregated results were supplied for the manuscript, internal consistency reliability using Cronbach's alpha was not calculated.

Actual economic outcomes, such as reduced post-harvest losses or increased farmer income, were not measured during the short testing period. In addition, complete security scanner outputs and manual penetration testing reports were not included in the results. Future evaluation should include raw response data, a longer field trial, transaction tracking, and formal security test documentation.

\subsection{Discussion Summary}

Overall, the results show that the Digital Farmers' Produce System was positively evaluated by its intended users. The system received strong ratings in usefulness, ease of use, attitude toward using, behavioral intention to use, functionality, and satisfaction. These findings indicate that the system has strong potential for adoption as a localized agricultural support platform for the Municipality of Sinacaban.

The pilot results also show that the system can help address important gaps in agricultural coordination, including limited produce visibility, weak farmer-buyer linkages, and lack of structured monitoring for the agriculture office. However, successful implementation will require user training, continued technical support, active participation from stakeholders, and future improvements based on actual field use.